%==============================================================================
% exhibits.tex
% Who Benefits from Food Tax Exemptions in Lower-Income Settings?
%
% Compiles every figure and table in the paper's order, with the paper's titles
% and notes, pointing each figure at the output file produced by this package
% (see docs/exhibit_manifest.md). Tables are reproduced inline from the paper.
%
% HOW TO BUILD
%   This file lives in docs/. Run pdflatex twice (it also builds from the
%   project root). Figures are read from Outputs/Figures via \graphicspath
%   below; edit that line if your Outputs folder sits elsewhere relative to
%   this file.
%
% MAPPING NOTES
%   - Figures 1 and 2 come from the R stage (figure_1.png, figure_2.png).
%   - Appendix A12/A13: the package filenames are swapped relative to the paper.
%       Paper Fig A12 (price-index trends) reads figure_a13.png
%       Paper Fig A13 (event study)        reads figure_a12.png
%   - External assets not produced by code (PNG Government notices Fig A8/A9,
%     the Stop&Shop photo Fig A16, and the Appendix C expert-survey PDF) are
%     shown as labeled placeholders so the auto-numbering stays exact. Swap in
%     the real image files where available.
%   - Expert figures A23 to A33 map to package figure_a23..a33 in the paper's
%     expert-1..expert-11 order.
%==============================================================================
\documentclass[12pt]{article}
\usepackage[margin=1in]{geometry}
\usepackage{graphicx}
\usepackage{float}
\usepackage{caption}
\usepackage{subcaption}
\usepackage{booktabs}
\usepackage{threeparttable}
\usepackage{array}
\usepackage{tabularx}
\usepackage{multirow}
\usepackage{multicol}
\usepackage{amsmath,amssymb}
\usepackage[table]{xcolor}
\usepackage{pdflscape}
\usepackage[hidelinks]{hyperref}

% Read figures from the package Outputs folder. Works whether this file sits in
% docs/ or at the project root, and whether figures are flat or in subfolders.
\graphicspath{%
  {Outputs/Figures/}{../Outputs/Figures/}%
  {Outputs/Figures/event_study/}{../Outputs/Figures/}%
  {Outputs/Figures/consumption_behavior/}{../Outputs/Figures/}}

% Placeholder for external assets not generated by this package
\newcommand{\externalasset}[1]{%
  \fbox{\parbox[c][4cm][c]{0.8\textwidth}{\centering\itshape
  External asset (not generated by this replication package):\\[4pt] #1}}}

\title{\vspace{-2cm}Exhibits: figures and tables\\[4pt]
\large Who Benefits from Food Tax Exemptions in Lower-Income Settings?\\
Evidence on Pass-Through and Incidence}
\date{}

\begin{document}
\maketitle
\thispagestyle{empty}
\vspace{-1cm}
\begin{center}\small Figures and tables in paper order. Figure titles and notes match the paper; each
figure is the file produced by the replication package (see \texttt{docs/exhibit\_manifest.md}).\end{center}
\clearpage

%==============================================================================
\section*{Main tables and figures}
%==============================================================================

% ---- Table 1: Summary of data sources (hand-built in the paper) -------------
\begin{table}[H]
    \centering
    \footnotesize
    \caption{Summary of data sources}
    \label{tab:datasets}
    \begin{threeparttable}
    \begin{tabularx}{\textwidth}{l>{\raggedright\arraybackslash}X>{\raggedright\arraybackslash}X>{\raggedright\arraybackslash}X>{\raggedright\arraybackslash}X}
        \toprule
         & \textbf{Supermarket census} & \textbf{Administrative data} & \textbf{Scraping of websites} & \textbf{Phone survey panel} \\
        \midrule
        Location    & Port Moresby & Port Moresby & Port Moresby & Nationwide \\[1mm]
        Unit        & Store & Store & Store & Household \\[1mm]
        Store type  & Formal & Formal & Formal & Formal + Informal \\[1mm]
        Coverage    & All & Subset & All online & Representative\\[1mm]
        Frequency   & Monthly & Monthly & Weekly & Monthly\\[1mm]
        Exempt items& 10 & 10 & 10 & 8 \\[1mm]
        Non-exempt items& 10 & 129 & 479 & 8 \\[1mm]
        Time period & May-Nov 25 & Dec 23 - May 26 & May-Aug 25 & Oct 23 - May 26 \\[1mm]
        Strengths   & All formal stores & Multi-year & All products & Consumption data \\[1mm]
        Limitations & Limited pre-trends & Non-representative & Only two chains & Self-reported \\
        \bottomrule
    \end{tabularx}
    \begin{tablenotes}[flushleft]
        \footnotesize
        \item \textit{Notes:} This table summarizes the four data sources used in the analysis. Location refers to geographic scope. Unit refers to whether the data covers stores or households. Store type refers to whether the data sources cover just formal or formal and informal stores. Coverage refers to whether the data source captures the universe of relevant outlets or a subset. Frequency refers to how often the data was collected. ``Exempt items'' and ``non-exempt items'' denote the number of treated and comparison-group product categories used in each source. Time period refers to the period for which this data source was collected. In the case of the scraped website data, there are 97 specific products within the broader 10 exempt food item categories.
    \end{tablenotes}
    \end{threeparttable}
\end{table}

% ---- Figure 1 ----------------------------------------------------------------
\begin{figure}[H]
    \centering
    \includegraphics[width=0.75\textwidth]{figure_1.png}
    \caption{Map of Port Moresby with supermarket locations}
    \caption*{\footnotesize \textit{Notes:} This figure maps the 83 formal supermarkets included in the Port Moresby supermarket census. Marker shading indicates stores' proximity to competitors, based on the Ahlfeldt--Redding--Sturm--Wolf proximity measure, with darker markers denoting greater local competitive intensity. High-proximity stores are concentrated in central Port Moresby.}
    \label{fig:sotres_location}
\end{figure}

% ---- Figure 2 ----------------------------------------------------------------
\begin{figure}[H]
    \centering
    \includegraphics[width=0.7\textwidth]{figure_2.png}
    \caption{Map of Papua New Guinea with respondents' locations}
    \caption*{\footnotesize \textit{Notes:} This figure maps the geographic distribution of households in the monthly phone survey panel. Shading indicates the number of unique respondent households observed in each district over the sample period. The figure illustrates the survey's broad national coverage.}
    \label{fig:png_map}
\end{figure}

% ---- Figure 3 ----------------------------------------------------------------
\begin{figure}[H]
    \centering
    \caption*{(a) Raw trends}
    \includegraphics[width=0.7\textwidth]{figure_3a.png}\\[4mm]
    \caption*{(b) Differences-in-Differences results}
    \includegraphics[width=0.7\textwidth]{figure_3b.png}
    \caption{Pass-through of VAT exemptions to prices in Port Moresby by data source}
    \caption*{\footnotesize \textit{Notes:} Panel (a) plots normalized price indices for VAT-exempt (treatment) and non-exempt (control) items by data source, with prices scaled to equal one in May 2025. Panel (b) reports difference-in-differences estimates of the price effect of the VAT exemption reform, relative to May 2025. Estimates are based on Equation 1. The vertical line marks the introduction of the VAT exemptions on 1 June 2025, and the horizontal reference line denotes full pass-through of a 10 percent VAT removal. Vertical bars show 95 percent confidence intervals.}
    \label{fig:combined_pass_through}
\end{figure}

% ---- Figure 4 (four panels) --------------------------------------------------
\begin{figure}[H]
    \centering
    \caption*{(a) Difference-in-Difference results by proximity}
    \includegraphics[width=0.8\textwidth]{figure_4a.png}\\[4mm]
    \caption*{(b) Distribution of pass-through rates}
    \includegraphics[width=0.8\textwidth]{figure_4b.png}
    \caption{Heterogeneity of pass-through based on level of competition}
    \label{fig:pass_through_by_competition}
\end{figure}
\begin{figure}[H]\ContinuedFloat
    \centering
    \caption*{(c) Differences in pass-through within the largest supermarket chain}
    \includegraphics[width=0.66\textwidth]{figure_4c.png}\\[4mm]
    \caption*{(d) Differences in pass-through online vs in-store}
    \includegraphics[width=0.62\textwidth]{figure_4d.png}
    \caption{Heterogeneity of pass-through based on level of competition (continued)}
    \caption*{\footnotesize \textit{Notes:} This figure examines heterogeneity in pass-through by local market competition using the supermarket census data. Panel (a) reports difference-in-differences estimates separately for stores in high-competition and low/mid-competition areas. Panel (b) shows the distribution of store-level pass-through rates three months after the reform. Panel (c) compares pass-through across stores within the largest supermarket chain, and Panel (d) compares online and in-store prices within the same chain. In Panels (a) and (d), the vertical line marks the reform date, and the dashed horizontal line indicates full pass-through of the VAT removal.}
\end{figure}

% ---- Table 2 -----------------------------------------------------------------
% Package output: 3_analysis/table2_passthrough.do -> table_2.tex
\begin{table}[H]
\centering
\caption{VAT pass-through estimates (p.p.)}
\label{tab:phone_pass_through_post_raw}
\begin{threeparttable}
\begin{tabular}{lccccc}
\toprule
Pass-through, p.p. & (1) & (2) & (3) & (4) & (5) \\
\midrule
Treatment & -3.1 & -0.9 & -0.0 & -1.0 & -1.0 \\
 & (3.1) & (2.6) & (3.8) & (2.5) & (2.5) \\
Treatment $\times$ Formal & --- & -4.9 & -7.3* & -1.3 & -1.2 \\
 &  & (3.9) & (4.3) & (3.3) & (3.3) \\
Treatment $\times$ Urban & --- & --- & --- & --- & 1.2 \\
 &  &  &  &  & (2.9) \\
Treatment $\times$ Formal $\times$ Urban & --- & --- & --- & --- & -6.7* \\
 &  &  &  &  & (3.9) \\
\midrule
Specification & Post & Post & Post & Post & Post \\
Coverage & All & All & Urban & Rural & All \\
$N$ & 129,921 & 129,921 & 43,725 & 86,196 & 129,921 \\
\bottomrule
\end{tabular}
\begin{tablenotes}[flushleft]
\scriptsize
\item \textit{Notes:} Coefficients are in p.p. of price change for VAT-exempt vs. control items. Prices are normalized to item-specific May 2025 means. Post is the simple average of post-reform event-study coefficients from June 2025 to May 2026. Full pass-through corresponds to $-9.1$ p.p. Formal is based on the store type where the purchase occurred. Conventional standard errors are clustered at the product level and include the standard finite-sample correction. Significance stars are based on wild-cluster bootstrap p-values. *** $p<0.01$, ** $p<0.05$, * $p<0.10$.
\end{tablenotes}
\end{threeparttable}
\end{table}

% ---- Figure 5 ----------------------------------------------------------------
\begin{figure}[H]
    \centering
    \caption*{(a) Share of foregone revenue accruing to households, by quintile}
    \includegraphics[width=0.8\textwidth]{figure5a.png}\\[4mm]
    \caption*{(b) Decomposition of consumer-side incidence}
    \includegraphics[width=0.9\textwidth]{figure5b.png}
    \caption{Incidence of VAT exemptions and decomposition of regressivity}
    \caption*{\footnotesize \textit{Notes:} Panel (a) reports the share of total foregone VAT revenue accruing to households in each expenditure quintile. The bars sum to 64 percent because the remaining share of foregone revenue is not passed through to households and is retained by stores and wholesalers. Panel (b) decomposes consumer-side incidence across quintiles. The blue series incorporates differences in consumption patterns only; the green series additionally incorporates heterogeneity in pass-through by store formality; and the red series further incorporates heterogeneity in pass-through by location and competition. In Panel (b), values are shown as deviations from equal incidence across quintiles, so the zero line corresponds to an equal per-quintile share of consumer benefits. Positive values indicate that a quintile captures more than an equal share of consumer benefits.}
    \label{fig:incidence}
\end{figure}

% ---- Table 3 -----------------------------------------------------------------
% Package output: 3_analysis/table3_alternative_policies.do -> table_3.tex
\begin{table}[H]
\centering
\caption{Cost of delivering USD 12 million to the poorest quintile under alternative policies}
\label{tab:counterfactual}
\begin{tabular}{lcc}
\hline\hline
Policy & Cost to deliver & Share of foregone revenue \\
\hline
Perfect targeting (Q1 only)  & USD 12 million  & 8\%   \\
Universal transfer (no targeting)     & USD 60 million & 40\%  \\
Universal transfer (50\% leakage) & USD 120 million   & 80\%  \\
Child grant and nutrition program & USD 150 million  & 100\% \\
VAT exemption (actual)       & USD 150 million  & 100\% \\
\hline\hline
\end{tabular}
\begin{minipage}{0.9\textwidth}
\vspace{0.5em}
\footnotesize
\textit{Notes:} This table compares the fiscal cost of alternative policies that would deliver the same monetary benefit to the poorest quintile as the VAT exemption reform. The benchmark assumes that the actual VAT exemption delivers approximately USD 12 million to the poorest quintile, or about 8 cents of each dollar of foregone revenue. The alternative rows report stylized transfer schemes under different targeting assumptions and implied leakage.
\end{minipage}
\end{table}

% ---- Figure 6 ----------------------------------------------------------------
\begin{figure}[H]
    \centering
    \includegraphics[width=0.75\textwidth]{figure6.png}
    \caption{Actual incidence vs. expert predictions of revenue lost across quintiles and stores}
    \caption*{\footnotesize \textit{Notes:} This figure compares the realized incidence of the VAT exemption reform with experts' ex ante predictions. Bars show the share of foregone revenue accruing to formal stores and to each household expenditure quintile. Realized incidence is computed from the phone survey using observed consumption patterns and estimated pass-through, while predicted incidence is based on the expert survey. The figure highlights the gap between expected and actual distributional outcomes.}
    \label{fig:incidence_actual_experts}
\end{figure}

\clearpage
%==============================================================================
\appendix
\setcounter{table}{0}
\setcounter{figure}{0}
\renewcommand\thefigure{A\arabic{figure}}
\renewcommand\thetable{A\arabic{table}}
\section*{Appendix A -- Additional Figures and Tables}
%==============================================================================

% A1
\begin{figure}[H]\centering
    \includegraphics[width=0.7\textwidth]{figure_a1.png}
    \caption{Average expenditure by item per quintile}
    \caption*{Note: This figure shows the average expenditure by item per quintile in Kina.}
    \label{fig:expenditure_item}
\end{figure}
% A2
\begin{figure}[H]\centering
    \includegraphics[width=0.7\textwidth]{figure_a2.png}
    \caption{Average expenditure by quintile per item}
    \caption*{Note: This figure shows the average expenditure by quintile per item in Kina.}
    \label{fig:expenditure_quintile}
\end{figure}
% A3
\begin{figure}[H]\centering
    \includegraphics[width=0.7\textwidth]{figure_a3.png}
    \caption{Average expenditure by quintile on basic items}
    \caption*{Note: This figure shows the average expenditure by quintile on basic items in Kina.}
    \label{fig:expenditure_basic}
\end{figure}
% A4
\begin{figure}[H]\centering
    \includegraphics[width=0.7\textwidth]{figure_a4.png}
    \caption{Average expenditure by quintile on quasi-luxury items}
    \caption*{Note: This figure shows the average expenditure by quintile on quasi-luxury items in Kina.}
    \label{fig:expenditure_quasilux}
\end{figure}
% A5
\begin{figure}[H]\centering
    \includegraphics[width=0.7\textwidth]{figure_a5.png}
    \caption{Quantity purchased by item per quintile}
    \caption*{Note: This figure shows the quantity purchased by item per quintile.}
    \label{fig:quantity_item}
\end{figure}
% A6
\begin{figure}[H]\centering
    \includegraphics[width=0.7\textwidth]{figure_a6.png}
    \caption{Quantity purchased by item by store type}
    \caption*{Note: This figure shows the quantity purchased by item by store type.}
    \label{fig:quantity_store}
\end{figure}
% A7
\begin{figure}[H]\centering
    \includegraphics[width=0.7\textwidth]{figure_a7.png}
    \caption{Type of consumption and expenditure by quintile}
    \caption*{Note: This figure shows the type of consumption and expenditure by quintile based on 2024 data from the phone survey.}
    \label{fig:type_cons_exp}
\end{figure}
% A8 (external: first PNG Government notice)
\begin{figure}[H]\centering
    \externalasset{Notice.png -- first official PNG Government notice}
    \caption{First official notice from the PNG Government about the VAT exemption reform}
    \caption*{Note: This figure shows the first official notice from the PNG Government about the VAT exemption reform.}
    \label{fig:notice}
\end{figure}

% Table A1: timeline (hand-built)
\begin{table}[H]\centering
\caption{Timeline of the VAT Exemptions Reform and Study Data Collection}
\label{tab:timeline_vat_exemptions}
\begin{tabular}{p{2.6cm} p{11.6cm}}
\hline\hline
\textbf{Date} & \textbf{Event} \\
\hline
Oct 2023 to May 2026 & Nationally representative, monthly \textbf{phone survey panel} (about 1,000 households/month) measuring quantities, expenditures, and store type (formal/informal) for 20 items.\\[2pt]
Dec 2023 to May 2026 & \textbf{Administrative data} collected monthly by (i) Central Bank of PNG research team and (ii) National Statistics Office CPI team; providing detailed price data for over 100 items from a subset of formal supermarkets in Port Moresby. \\[2pt]
29 Nov 2024 & PNG government first publicly stated the \emph{possibility} of removing VAT from some basic items in mid-2025; details were limited and credibility unclear. \\[2pt]
19 May 2025 & Policymakers announced that VAT exemptions for 10 food items (and 3 non-food items) would take effect on \textbf{June 1, 2025} (i.e., 12 days later). \\[2pt]
May-Aug 2025 & Weekly web-scraped all \textbf{online prices} from the two chains offering online shopping in Port Moresby (including the use of ``time machine'' software to check prices before reform announced).\\[2pt]
May-Nov 2025 & Monthly in-store \textbf{supermarket census} collected prices for 20 food items (10 exempt and 10 non-exempt) in all 83 formal supermarkets in Port Moresby.  \\[2pt]
1 June 2025 & Reform implemented: VAT on 10 food items set to zero-rated (functionally close to an exemption in practice). Exemptions scheduled to remain in place for at least 13 months. \\[2pt]
July-Sep 2025 & Expert prediction survey hosted on SSPP to benchmark beliefs about pass-through and incidence against realized outcomes. \\[2pt]
8 Aug 2025 & Policymakers issued revised guidance clarifying the details of the reform and outlining plans to monitor its impact on prices. \\[2pt]
25 Nov 2025 & PNG government announced an extension of the VAT exemption policy to the end of 2026 (i.e., 19 months in total). \\[2pt]
Dec 2025 & Phone survey question about how familiar households are with the details of the reform six months after implementation. \\
\hline\hline
\end{tabular}
\begin{flushleft}
\footnotesize \textit{Notes:} The table lists both policy milestones and the data-collection windows used to estimate pass-through and incidence.
\end{flushleft}
\end{table}

% A9 (external: second PNG Government notice, 3 pages)
\begin{figure}[H]\centering
    \externalasset{Figure\_1\_Page\_1 / \_2 / \_3 -- second official PNG Government notice (3 pages)}
    \caption{Second official notice from the PNG Government about the VAT exemption reform}
    \caption*{Note: This figure shows the second official notice from the PNG Government about the VAT exemption reform.}
    \label{fig:notice2}
\end{figure}

% Table A2: representativeness (package: survey_census_representativeness.tex; numbers reproduced)
\begin{table}[H]\centering
  \begin{tabular}{lcccc} \toprule & \multicolumn{3}{c}{Survey} & Census \\ \cmidrule(lr){2-4} \cmidrule(lr){5-5} & All months & 1-month & Monthly average &  \\  \midrule
BOUGAINVILLE & 4.81 & 5.76 & 4.62 & 3.51 \\
CENTRAL & 3.00 & 2.49 & 2.73 & 3.01 \\
CHIMBU & 5.52 & 7.20 & 5.73 & 5.54 \\
EAST NEW BRITAIN & 4.93 & 5.10 & 4.51 & 4.25 \\
EAST SEPIK & 5.84 & 6.02 & 5.91 & 6.36 \\
EASTERN HIGHLANDS & 8.68 & 7.72 & 7.72 & 9.97 \\
ENGA & 4.65 & 5.63 & 5.50 & 5.56 \\
GULF & 2.01 & 1.96 & 2.45 & 1.88 \\
HELA & 3.23 & 5.10 & 4.25 & 4.75 \\
JIWAKA & 4.62 & 4.97 & 5.33 & 4.74 \\
MADANG & 6.07 & 5.24 & 6.24 & 6.27 \\
MANUS & 2.49 & 2.62 & 2.52 & 0.75 \\
MILNE BAY & 2.37 & 1.83 & 2.10 & 4.02 \\
MOROBE & 8.28 & 6.54 & 7.62 & 9.46 \\
NATIONAL CAPITAL DISTRICT & 5.68 & 5.89 & 4.15 & 3.46 \\
NEW IRELAND & 3.43 & 2.62 & 3.45 & 2.16 \\
NORTHERN (ORO) & 2.09 & 2.09 & 2.55 & 2.48 \\
SOUTHERN HIGHLANDS & 4.93 & 6.54 & 6.34 & 6.40 \\
WEST NEW BRITAIN & 4.50 & 2.88 & 3.69 & 3.69 \\
WEST SEPIK (SANDAUN) & 4.54 & 3.66 & 4.07 & 3.27 \\
WESTERN & 1.78 & 1.05 & 1.54 & 2.28 \\
WESTERN HIGHLANDS & 6.55 & 6.81 & 6.98 & 6.19 \\
\bottomrule \end{tabular}
    \caption{Distribution of survey respondents relative to population}
    \caption*{Note: This table shows the distribution of survey respondents relative to the population.}
    \label{tab:rep}
\end{table}

% Table A3: familiarity (package: table_a3.tex)
\begin{table}[H]
\centering
\caption{Familiarity with National Government's Efforts to Reduce Tax}
\label{tab:knowledge}
\begin{tabular}{lrrr}
\hline\hline
Response & Freq. & Percent & Cum. \\
\hline
Not at all    & 510   & 38.06 & 38.06  \\
A little      & 536   & 40.00 & 78.06  \\
Somewhat      & 248   & 18.51 & 96.57  \\
Very Familiar &  46   &  3.43 & 100.00 \\
\hline
Total         & 1,340 & 100.00 &       \\
\hline\hline
\end{tabular}
\end{table}

% A10
\begin{figure}[H]\centering
    \includegraphics[width=0.75\textwidth]{figure_a10.png}
    \caption{Average monthly expenditure on exempt items}
    \caption*{Note: This figure shows the average monthly expenditure on exempt items.}
    \label{fig:exempt}
\end{figure}
% A11
\begin{figure}[H]\centering
    \includegraphics[width=0.75\textwidth]{figure_a11.png}
    \caption{Average monthly expenditure on non-exempt items}
    \caption*{Note: This figure shows the average monthly expenditure on non-exempt items.}
    \label{fig:non_exempt}
\end{figure}
% A12 (paper: price-index trends -> package figure_a13.png)
\begin{figure}[H]\centering
    \includegraphics[width=0.8\textwidth]{figure_a13.png}
    \caption{Raw trends in prices over the last 14 years according to NSO administrative data}
    \caption*{Note: This figure shows the raw trends in prices over the last 14 years according to NSO administrative data.}
    \label{fig:raw_long}
\end{figure}
% A13 (paper: event study -> package figure_a12.png)
\begin{figure}[H]\centering
    \includegraphics[width=0.8\textwidth]{figure_a12.png}
    \caption{Differences-in-Differences results with different comparison groups over the last four years according to NSO administrative data}
    \caption*{Note: This figure shows the Differences-in-Differences results with different comparison groups over the last four years according to NSO administrative data.}
    \label{fig:did_long}
\end{figure}
% A14
\begin{figure}[H]\centering
    \includegraphics[width=0.75\textwidth]{figure_a14.png}
    \caption{Raw trends of prices in Port Moresby by data source}
    \caption*{Note: This figure shows the raw trends results by data source.}
    \label{fig:a14}
\end{figure}
% A15
\begin{figure}[H]\centering
    \includegraphics[width=0.75\textwidth]{figure_a15.png}
    \caption{Pass-through of VAT exemptions to prices in Port Moresby by data source}
    \caption*{Note: This figure shows the difference-in-differences results by data source.}
    \label{fig:a15}
\end{figure}
% A16 (external: Stop&Shop advertisement photo)
\begin{figure}[H]\centering
    \externalasset{shop.jpg -- Stop and Shop advertisement}
    \caption{Advertisement from Stop and Shop}
    \caption*{\footnotesize Note: This figure presents an advertisement from the largest supermarket chain, Stop and Shop.}
    \label{fig:stopandshop}
\end{figure}
% A17
\begin{figure}[H]\centering
    \includegraphics[width=0.8\textwidth]{figure_a17.png}
    \caption{Results of pass-through question in business sentiment survey}
    \caption*{\footnotesize Note: This figure presents the results of the pass-through question in the business sentiment survey administered by the central bank.}
    \label{fig:pass}
\end{figure}
% A18
\begin{figure}[H]\centering
    \includegraphics[width=0.8\textwidth]{figure_a18.png}
    \caption{Results of competition question in business sentiment survey}
    \caption*{\footnotesize Note: This figure presents the results of the competition question in the business sentiment survey administered by the central bank.}
    \label{fig:comp}
\end{figure}
% A19 (four panels: package figure_a19a..d)
\begin{figure}[H]\centering
    \begin{minipage}[t]{0.48\textwidth}\centering
        \caption*{(a) Pass-through rate nationwide}
        \includegraphics[width=\textwidth]{figure_a19a.png}
    \end{minipage}\hfill
    \begin{minipage}[t]{0.48\textwidth}\centering
        \caption*{(b) By store formality}
        \includegraphics[width=\textwidth]{figure_a19b.png}
    \end{minipage}\\[3mm]
    \begin{minipage}[t]{0.48\textwidth}\centering
        \caption*{(c) Urban areas}
        \includegraphics[width=\textwidth]{figure_a19c.png}
    \end{minipage}\hfill
    \begin{minipage}[t]{0.48\textwidth}\centering
        \caption*{(d) Rural areas}
        \includegraphics[width=\textwidth]{figure_a19d.png}
    \end{minipage}
    \caption{Pass-through nationwide based on the phone survey data}
    \caption*{\footnotesize \textit{Notes:} This figure reports difference-in-differences estimates using the nationally representative phone survey. Panel (a) shows the average pass-through nationwide. Panel (b) disaggregates by store formality. Panels (c) and (d) further split the results by urban and rural location. The vertical line marks the June 2025 reform, and the dashed horizontal line indicates full pass-through of the VAT removal.}
    \label{fig:pass_through_heterogeneity}
\end{figure}
% A20
\begin{figure}[H]\centering
    \includegraphics[width=0.7\textwidth]{figure_a20.png}
    \caption{Raw trends in prices in treatment and control groups in the phone survey data}
    \caption*{Note: This figure shows the raw trends in prices in treatment and control groups in the phone survey data.}
    \label{fig:raw_tvsc}
\end{figure}
% A21
\begin{figure}[H]\centering
    \includegraphics[width=0.7\textwidth]{figure_a21.png}
    \caption{Raw trends in prices in treatment group between formal and informal stores in the phone survey data}
    \caption*{Note: This figure shows the raw trends in prices in treatment group between formal and informal stores in the phone survey data.}
    \label{fig:raw_fvsi}
\end{figure}
% A22
\begin{figure}[H]\centering
    \includegraphics[width=0.8\textwidth]{figure_a22.png}
    \caption{Preferences for alternative policies}
    \caption*{\footnotesize Note: This figure shows the results of a question about preferences for alternative policies asked in the April 2026 round of the phone survey. The exact wording was as follows (and respondents selected their first and second preference): ``Imagine the government was considering introducing a program to support poor and middle-class families. In your view, which of the following would be the best way to support poor and middle-class families? provide a cash transfer for all families with children under the age of five; remove taxes from basic food items sold in supermarkets; remove taxes from petrol and other fuels; provide free textbooks for schools; provide a small amount of free phone credit for all phone users; reduce the level of government debt.''}
    \label{fig:pref}
\end{figure}

% Table A4: HTE (package: table_a4.tex)
\begin{table}[H]
\centering
\caption{Treatment effects on support for removing the tax on food items, by prior policy preference}
\label{tab:hte_foodtax}
\footnotesize
\begin{tabular}{lc r@{\hspace{3pt}}l r@{\hspace{3pt}}l}
\toprule
 & & \multicolumn{2}{c}{Unweighted} & \multicolumn{2}{c}{Weighted} \\
\cmidrule(lr){3-4}\cmidrule(lr){5-6}
Subgroup & $N$ & \multicolumn{2}{c}{(robust SE)} & \multicolumn{2}{c}{(robust SE)} \\
\midrule
Overall (ATE) & 995 & $-0.119^{*}$ & $(0.063)$ & $-0.106$ & $(0.091)$ \\
\addlinespace
\multicolumn{6}{l}{\textit{Preferred way to support households (April 2026):}}\\
\quad Prefers removing food tax & 340 & $+0.049$ & $(0.113)$ & $+0.186$ & $(0.151)$ \\
\quad Prefers another policy & 412 & $-0.230^{**}$ & $(0.090)$ & $-0.285^{**}$ & $(0.119)$ \\
\addlinespace
\multicolumn{6}{l}{\quad\textit{Another policy, by type:}}\\
\quad\quad Cash transfer (under-5s) & 86 & $-0.796^{***}$ & $(0.200)$ & $-0.914^{***}$ & $(0.288)$ \\
\quad\quad Reduce government debt & 152 & $-0.036$ & $(0.151)$ & $-0.207$ & $(0.213)$ \\
\quad\quad Free textbooks & 109 & $-0.285^{*}$ & $(0.173)$ & $-0.199$ & $(0.197)$ \\
\quad\quad Remove tax on fuel & 46 & $+0.200$ & $(0.262)$ & $-0.019$ & $(0.242)$ \\
\quad\quad Free phone credit & 19 & \multicolumn{2}{c}{--} & \multicolumn{2}{c}{--} \\
\addlinespace
\multicolumn{6}{l}{\textit{Interaction tests:}}\\
\quad Prefers food tax $\times$ Treat & & $+0.286^{**}$ & $(0.143)$ & $+0.467^{**}$ & $(0.192)$ \\
\quad Prefers cash $\times$ Treat & & $-0.775^{***}$ & $(0.212)$ & $-0.927^{***}$ & $(0.311)$ \\
\quad Joint test, all groups ($F$) & & \multicolumn{2}{c}{$p=0.007$} & \multicolumn{2}{c}{$p=0.049$} \\
\bottomrule
\end{tabular}
\begin{minipage}{0.97\textwidth}
\vspace{4pt}\scriptsize
\textit{Notes:} The outcome is support for removing the tax on selected food items, on a 1--5 scale (5 = strongly support; control-group mean 3.76). Cells report subgroup average treatment effects, conditioning on the orthogonal question-version indicator, with heteroskedasticity-robust (HC1) standard errors in parentheses. The unweighted estimates (preferred specification) are shown alongside weighted estimates that apply the household survey weights as a robustness check. Treatment was randomly assigned in the May 2026 round ($N=995$); respondents' preferred way to support households is taken from the same households' April 2026 responses, which restricts the heterogeneity analysis to the 752 households observed in both rounds. The interaction rows report the treatment $\times$ moderator coefficient (food tax versus all other preferences; cash transfer versus all others); the final row reports the $p$-value of a joint $F$-test across the six preference groups. The smallest group (free phone credit) is too small to estimate separately but is retained in the joint test. $^{*}p<0.10$, $^{**}p<0.05$, $^{***}p<0.01$.
\end{minipage}
\end{table}

% Table A5: expert background (package: table_a5.tex)
\begin{table}[H]\centering
  \begin{tabular}{lc} \midrule
   \noalign{\smallskip} Finished (\%) & 73.84 \\
   Median duration (minutes) & 8.97 \\
   Researcher (\%) & 62.44 \\
   Economics (\%) & 92.44 \\
   \shortstack[l]{Research experience in: \\ \quad{Public Finance (\%)}} & 52.54 \\
   \quad{Tax Incidence (\%)} & 31.07 \\
   \quad{Development (\%)} & 58.76 \\
   \quad{Industrial Organization (\%)} & 16.38 \\
   \quad{Low-middle income countries (\%)} & 62.15 \\
   \midrule Number of responses & 237  \\ Number of complete responses & 175 \\ \bottomrule \end{tabular}
    \caption{Background information about expert prediction survey}
    \caption*{Note: This table shows background information about the expert prediction survey.}
    \label{tab:expert}
\end{table}

% A23 - A33 expert figures
\begin{figure}[H]\centering
    \includegraphics[width=0.8\textwidth]{figure_a23.png}
    \caption{Experts' predictions of pass-through by data source}
    \caption*{Note: This figure shows experts' predictions of pass-through by data source.}
    \label{fig:expert1}
\end{figure}
\begin{figure}[H]\centering
    \includegraphics[width=0.8\textwidth]{figure_a24.png}
    \caption{Experts' predictions of variation in pass-through}
    \caption*{Note: This figure shows experts' predictions of variation in pass-through.}
    \label{fig:expert2}
\end{figure}
\begin{figure}[H]\centering
    \includegraphics[width=0.8\textwidth]{figure_a25.png}
    \caption{Experts' predictions of items with high pass-through}
    \caption*{Note: This figure shows experts' predictions of items with high pass-through.}
    \label{fig:expert3}
\end{figure}
\begin{figure}[H]\centering
    \includegraphics[width=0.8\textwidth]{figure_a26.png}
    \caption{Experts' predictions of items with low pass-through}
    \caption*{Note: This figure shows experts' predictions of items with low pass-through.}
    \label{fig:expert4}
\end{figure}
\begin{figure}[H]\centering
    \includegraphics[width=0.8\textwidth]{figure_a27.png}
    \caption{Experts' predictions of pass-through by store type}
    \caption*{Note: This figure shows experts' predictions of pass-through by store type.}
    \label{fig:expert5}
\end{figure}
\begin{figure}[H]\centering
    \includegraphics[width=0.8\textwidth]{figure_a28.png}
    \caption{Experts' predictions of changes in pass-through over time}
    \caption*{Note: This figure shows experts' predictions of changes in pass-through over time.}
    \label{fig:expert6}
\end{figure}
\begin{figure}[H]\centering
    \includegraphics[width=0.8\textwidth]{figure_a29.png}
    \caption{Experts' predictions of size of pass-through}
    \caption*{Note: This figure shows experts' predictions of size of pass-through.}
    \label{fig:expert7}
\end{figure}
\begin{figure}[H]\centering
    \includegraphics[width=0.8\textwidth]{figure_a30.png}
    \caption{Experts' reasons for substantial pass-through}
    \caption*{Note: This figure shows experts' reasons for substantial pass-through.}
    \label{fig:expert8}
\end{figure}
\begin{figure}[H]\centering
    \includegraphics[width=0.8\textwidth]{figure_a31.png}
    \caption{Experts' reasons for small pass-through}
    \caption*{Note: This figure shows experts' reasons for small pass-through.}
    \label{fig:expert9}
\end{figure}
\begin{figure}[H]\centering
    \includegraphics[width=0.8\textwidth]{figure_a32.png}
    \caption{Experts' predictions of incidence of reform by quintile}
    \caption*{Note: This figure shows experts' predictions of incidence of reform by quintile.}
    \label{fig:expert10}
\end{figure}
\begin{figure}[H]\centering
    \includegraphics[width=0.8\textwidth]{figure_a33.png}
    \caption{Experts' predictions of incidence of reform by quintile (boxplot)}
    \caption*{Note: This figure shows experts' predictions of incidence of reform by quintile (boxplot).}
    \label{fig:expert11}
\end{figure}
% A34, A35 expert heterogeneity
\begin{figure}[H]\centering
    \includegraphics[width=\linewidth]{figure_a34.png}
    \caption{Heterogeneity in experts' predictions of pass-through}
    \caption*{Note: This figure shows heterogeneity in experts' pass-through predictions.}
    \label{fig:expert12}
\end{figure}
\begin{figure}[H]\centering
    \includegraphics[width=\linewidth]{figure_a35.png}
    \caption{Heterogeneity in experts' predictions of incidence}
    \caption*{Note: This figure shows heterogeneity in experts' predictions of incidence.}
    \label{fig:expert13}
\end{figure}
% A36
\begin{figure}[H]\centering
    \includegraphics[width=0.8\textwidth]{figure_a36.png}
    \caption{Differences-in-Differences results by chain vs independent stores}
    \caption*{\footnotesize Note: This figure presents event study estimates including the average of post-reform event study coefficients (reference = May 2025), by chain vs independent stores (balanced panel).}
    \label{fig:aa2}
\end{figure}
% A37
\begin{figure}[H]\centering
    \includegraphics[width=0.8\textwidth]{figure_a37.png}
    \caption{Heterogeneity of pass-through based on chain vs independent stores}
    \caption*{Note: This figure shows how the pass-through rate varies by chain vs independent stores, using the supermarket census data.}
    \label{fig:aa3}
\end{figure}
% A38
\begin{figure}[H]\centering
    \includegraphics[width=0.75\textwidth]{figure_a38.png}
    \caption{Differences-in-Differences results by product type}
    \caption*{\footnotesize Note: This figure presents event study estimates including the average of post-reform event study coefficients (reference = May 2025), by basic vs luxury products (balanced panel).}
    \label{fig:aa4}
\end{figure}

\clearpage
\section*{Appendix C -- Expert Prediction Survey}
\begin{figure}[H]\centering
    \externalasset{Expertsurvey.pdf -- full expert prediction survey instrument}
\end{figure}

\end{document}
